%==============================================================================
\section{Experimental Design}
\label{sec:experimental-design}
%==============================================================================

\textbf{Datasets.} We use two public benchmarks. \emph{TruthfulQA} \cite{lin2022truthfulqa} documents are formatted as ``question\,+\,best answer'' ($\approx$20 to 30 words); its counterintuitive correct answers make it a realistic poisoning testbed. \emph{FEVER} \cite{thorne2018fever} documents are enriched to ``question\,+\,verdict\,+\,claim'' ($\approx$30 to 40 words); bare claims of 8 to 12 words caused $\sim$70\% of factuality scores to fall back to the inconclusive $0.5$ baseline, motivating the enrichment. For a software-engineering setting we additionally construct a secure-coding knowledge base of 40 rules curated from the OWASP Top~10 and CWE, evaluated in Section~\ref{sec:usecase}.

\textbf{Protocol.} Each run has two parts: 50 clean queries (Part~A) and the same 50 queries against a corpus with 30\% poisoned documents (Part~B), for 100 samples per run. We evaluate four rule-based poisoning strategies: \emph{contradiction} (appends a contradicting statement), \emph{instruction injection} (appends an override directive), \emph{entity swap} (in-place replacement of entities/numbers, with no appended artifacts), and \emph{subtle manipulation} (false qualifiers/hedging). A \emph{mixed} setting assigns strategies randomly across poisoned documents. Part~A contributes 50 clean samples and Part~B contributes 50 samples of which 30\% ($\approx$15) are labelled poisoned, so a 100-sample run contains about 15 poisoned and 85 clean samples. A naive always-trust baseline that labels every sample clean is therefore correct on 85 clean samples, i.e.\ $\approx$85\% accuracy; hence $\Delta$ and F1, not accuracy, are the informative metrics.

\textbf{Models and metrics.} Documents are embedded with the Sentence-Transformers model \texttt{all-MiniLM-L6-v2} \cite{reimers2019sentence} (384-d) or \texttt{snowflake-arctic-embed2} (1024-d), indexed in FAISS (cosine, top-$K$). Generation uses \texttt{llama3.3:70b} (primary) or \texttt{qwen3.5:35b}; NLI runs on CPU. We report accuracy, precision, recall, F1, and \emph{trust score separation} $\Delta = \bar{T}_{\text{clean}} - \bar{T}_{\text{poisoned}}$, a threshold-independent measure of discriminability. The comparison point is a \emph{naive always-trust} RAG that never flags; under 30\% poisoning its accuracy ($\approx$85\%) reflects class imbalance, not detection, so $\Delta$ and F1 are the informative metrics. We add 95\% confidence intervals: Wilson intervals for proportions and a percentile bootstrap ($B{=}20\,000$) for F1.

\textbf{Ground truth.} The clean and poisoned corpora are paired by index, so each query maps to a fixed source document. A sample is labelled \emph{poisoned} when its aligned source document was modified by a poisoning strategy in Part~B, and \emph{clean} otherwise. For example, if the document answering ``\emph{How should passwords be stored?}'' was poisoned, that query's sample is labelled poisoned even when the retriever also returns clean neighbours. This document-level label fits the paired benchmark but is stricter than a retrieval-grounded label, which would count a sample as poisoned only when a poisoned document enters the top-$K$ set.

\textbf{Reproducibility.} Generation uses temperature $0.7$, \texttt{max\_tokens}$=512$, and provider-default top-$p$; NLI (\texttt{facebook/bart-large-mnli}) runs on CPU; retrieval uses FAISS (cosine) over $512$-token chunks (overlap $50$) at $K{=}5$. Sampling and poison assignment use a fixed seed ($42$), while LLM decoding stays stochastic (quantified by the variance in Section~\ref{sec:stability}). Calibration fits $\tau$ on a held-out clean split (10th percentile) evaluated on disjoint clean/poisoned samples. All LLMs run on the FARMI/Ollama endpoint; code, prompts, and the attack generator are released.
